\section{Introduction}
\label{sec:intro}

The rapid adoption of large language models (LLMs) in production systems
has generated a body of prompt engineering heuristics: add role descriptions,
include worked examples, append step-by-step instructions, and specify output
formats in detail~\citep{brown2020gpt3,wei2022cot,zhou2022large}.
Official guidance from major providers recommends structured prompts with
role, objective, instructions, reasoning steps, output format, examples,
and context, each section adding tokens~\citep{sahoo2024systematic,schulhoff2024prompt}.
These practices implicitly assume a monotone relationship between prompt
elaboration and response quality.
But is this assumption empirically justified, and if so, \emph{up to what
point}?

Token costs are not merely economic.
Longer prompts consume more inference compute, fill context windows faster,
and can \emph{degrade} performance when key information is buried in
irrelevant context~\citep{liu2023lost}.
The prompt compression literature~\citep{jiang2023llmlingua,pan2024llmlingua2}
has demonstrated that many tokens are expendable without quality loss, but
these methods operate on already-verbose prompts and assume compression is
safe.
A more fundamental question is: \emph{when does prompt elaboration stop
helping in the first place?}

We address this through a controlled \textbf{prompt sensitivity study}.
For each of six tasks we construct seven strictly additive prompt templates:
each template is a proper superset of the previous, adding exactly one
\emph{mechanism} of elaboration (task label, output format, field definitions,
persona, guidelines, worked example).
Because level $k{+}1$ contains every token of level $k$ plus one new mechanism,
any quality difference between adjacent levels isolates that mechanism's
contribution.
We run every (model, task, level, example) combination for seven LLMs and
measure where quality plateaus.
Crucially, for the four tasks with objective ground truth (classification,
product extraction, QA, math) we score quality with \emph{deterministic}
metrics as the primary measure, reserving the LLM judge for tasks where no
clean heuristic exists; this avoids resting the conclusions on the reliability
of a single LLM evaluator.
We further disentangle information content from token count through two
complementary controls: a randomised layer-ordering ablation that shuffles
which mechanism appears at each token count, and an irrelevant-padding
experiment that matches token counts with meaningless filler.

Our central result is that, for structured tasks, quality gain
\textbf{concentrates at a single schema-resolving layer}: the layer that
communicates the output format or label space the model cannot infer from the
input alone.
On the deterministic metric, this layer captures the overwhelming majority of
the total L1$\to$L7 gain:
\begin{itemize}
  \item \textbf{Classification:} the task label (L1$\to$L2) captures
    \textbf{87\%} of total gain.
  \item \textbf{Product extraction:} the JSON format specification
    (L2$\to$L3) captures \textbf{92\%} of total gain, with quality rising from
    near-zero to ${>}0.90$; 7/7 models produce significant saturating fits
    ($R^2 > 0.99$).
  \item \textbf{Named-entity recognition (held-out):} the format specification
    (L2$\to$L3) captures \textbf{95\%} of total gain (4/5 significant),
    confirming the mechanism on a task chosen \emph{after} the main study.
\end{itemize}
Everything else (definitions, personas, guidelines, even the worked example)
adds little once the schema layer is present.
The LLM judge had attributed product-extraction gain partly to the worked
example, suggesting a ``distributed'' pattern; the ground-truth metric reveals
this as a judge artefact and shows the gain is concentrated at the format
layer (Section~\ref{sec:result-judge}).
A hard-example stress test and the held-out task then sharpen the full picture
into three mechanism-anchored classes:
\textbf{schema-gated} (classification, extraction, NER, where one layer unlocks
the task), \textbf{difficulty-gated} (instruction following, where sensitivity
grows $6.6\times$ with harder, multi-constraint examples), and
\textbf{knowledge-gated} (QA, math, summarisation, at ceiling from the bare
input because the answer depends on parametric knowledge, not prompt form).
Two controls confirm the mechanism is information, not length: reordering the
layers moves the token threshold but preserves the saturation shape, and
matched-token irrelevant filler produces flat or declining quality (1/26
conditions significant).

\paragraph{Contributions.}
\begin{enumerate}
  \item A reusable \textbf{additive prompt-level methodology}, with randomised
    ablation and padding controls and \emph{ground-truth-primary} evaluation,
    for measuring which prompt mechanisms drive quality, with open-source code
    and data.
  \item The \textbf{schema-gating finding}: for structured tasks, a single
    schema-resolving layer captures 87--95\% of all quality gain, measured on
    deterministic metrics (7/7 significant for product extraction,
    $R^2 > 0.99$) and confirmed on a \textbf{held-out task} (NER), making the
    taxonomy predictive rather than merely descriptive.
  \item Evidence that saturation is driven by \textbf{information
    sufficiency}, not token count, via a randomised layer-ordering ablation
    (2,986 runs) and an irrelevant-padding control (1,470 runs).
  \item A \textbf{difficulty-resolved} account of apparent insensitivity: a
    hard-example stress test on GSM8K and HotpotQA shows QA is genuinely
    knowledge-gated, while instruction-following sensitivity grows $6.6\times$
    with harder examples.
  \item \textbf{Reliability and stability evidence}: deterministic scoring of
    all ground-truth tasks, an independent second judge (3,915 re-evaluations),
    and repeated runs (effect-to-noise 20--139) show the effects far exceed
    both evaluator disagreement and sampling noise.
\end{enumerate}

The total experimental scope is 21,858 evaluations across the main study,
replication, ablation, padding control, second-judge validation, held-out NER,
hard-example stress test, sampling-stability repeats, and sub-8B experiments.
All prompts, evaluation code, and data are included in the supplementary
material.
